\subsection{6. 讨论}\label{discussion}

前几节给出的形式化模型与实现，引入了一种用于动态可组合性的编程范式。本节考察该范式如何扩展到更广泛的工程关切，并讨论其中存在的设计张力与开放问题。

\subsection{6.1. 系统边界}\label{system-boundary}

第~3.1 节的每一个效应都带有一个逆，而该逆的具体含义由系统边界决定。边界把系统所运行于其中的环境划分为两个部分。(1) 当系统能够独占性地修改某个位置、并能恢复该位置在此次修改之前的状态时，这个位置就位于边界之内，因此对它的操作会被 $\Gamma$ 跟踪，并可在日后恢复。(2) 当上述任一能力缺失时，这个位置位于边界之外，对它的操作表现为 \(\mathrm { i d } _ { \Gamma }\)，因而既不被跟踪、也不被恢复。本节展开这个边界的性质，以及这些性质对恢复的后果。

来自余效应的边界。余效应通过把外部位置具体化来移动边界：它把对该位置的每一次访问都限定在一组由它提供的操作之内，而它能为其间每个操作都提供逆，于是原本表现为 \(\mathrm { i d } _ { \Gamma }\) 的操作，转而进入 $\Gamma$ 被跟踪并得以恢复。因此，边界是按位置而非按介质来划定的，因为上述两种能力都是某个位置的属性，而具体化改变的是访问该位置的方式，其介质则保持不变。例如，一块内存区域当仅由系统写入时位于边界之内，当其他进程也写入它时则位于边界之外；一个文件当只有系统能够触达它时——如私有路径下的临时文件——位于边界之内，当它是其他程序可读写的路径时则位于边界之外。移动边界本身是一种权衡，权衡的双方是：环境是否为某个位置提供可逆语义，以及提供这些语义在每次访问时付出多少代价。我们把这一权衡所暗示的协同设计留待第~6.7 节讨论。

67

获取与发出。到达边界之外的操作通常分两个阶段进行。(1) 在获取阶段，操作获得访问权，并在边界之内安装一条记录：open 安装一个由 close 移除的描述符，malloc 保留一块由 free 释放的内存，fork 启动一个由 kill 终止的子进程。这条记录本身是具体化该位置的余效应的一部分，例如它维护的映射中的一个条目，而安装这个条目是一个可逆效应。这条记录同时也是数据得以离开的通道。(2) 在发出阶段，操作通过该通道推送数据，就像 write 交给文件的字节、或 send 放到线路上的数据报那样，而推送表现为 \(\mathrm { i d } _ { \Gamma }\)，把数据留在其他方可以读写的位置。因此，这两个阶段落在边界的两侧：获取停留在边界之内，而发出则跨越到边界之外。

扣留与补偿。一个仍必须从发出中恢复的系统，有两种可用的方法。其一是扣留发出，直到产生它的状态确定能够持久化为止，这正是回滚恢复中的输出提交问题 {[}48{]}。其二是补偿 {[}49{]}：用一种由应用程序提供的等价关系把状态恢复到某一程度，该关系比定义~33 中的 ≃ 更粗，例如删除一个刚创建的文件，或退还一笔已收取的费用。这类操作按与逆相同的 LIFO 顺序复合，因此第~3.1 节的复合性质可以转移到它们之上。元理论则不然：定义~60 的交换性是针对 ≃ 证明的，必须针对更粗的那个关系重新建立。

\subsection{6.2. 服务多路复用}\label{service-multiplexing}

OSGi {[}50{]} 等动态组件平台围绕服务来组织组合：服务是功能单元，由提供者在某个接口下发布，由消费者绑定。Cordis 的余效应模型呼应了这一概念，服务对应于键背后的接口。提供某个服务的组件是它的提供者，注入某个服务的组件是它的消费者。一个服务可能由多个提供者实现，而这种多重性有两种实现形式。(1) 独占绑定：多个实现共享一个接口，但同一时刻至多绑定一个；编排器选择绑定哪一个实现，而在它们之间切换需要卸载一个提供者并加载另一个，这会在瞬间扰动每个消费者的依赖。(2) 服务代理：一个作为该接口入口的中央服务，被底层提供者与消费者共同注入，从而使多个提供者并存，由代理在各提供者之间分派每个请求。与独占绑定相比，代理吸收了这种扰动：更新一个底层提供者时，代理保持原位，于是消费者看到自己的依赖没有任何变化，也不会触发重载。

服务代理支撑三种能力：负载均衡、滚动更新与跨进程调用。

负载均衡。当多个提供者并存时，代理按照可配置的策略（例如轮询、最少负载、按延迟加权）或消费者指定的显式目标，在各提供者之间分发请求。由于提供者是普通组件，可以添加或移除它们来扩大或缩小容量；每个提供者通过一个可逆效应向代理注册，因此卸载它会使注册被撤销，并自动从代理的路由集合中移除。

68

滚动更新。在运行时升级某个服务实现，可归结为一次受控的提供者迁移 {[}51, 52{]}。为完成这次迁移，新的提供者作为一条额外的纤维被加载并向代理注册；一旦它进入 ACTIVE 状态，流量便逐渐从旧提供者转向新提供者（例如通过调整选择权重），而旧提供者在其不再承载在途请求后被卸载。这种提供者迁移，把传统上属于基础设施层面的操作（例如容器编排、蓝绿部署）转变为应用层面的组合模式。

跨进程调用。服务代理也可以跨进程边界使用 {[}53{]}。每个进程承载自己的 Cordis 上下文，其中带有本地提供者；一个协调组件把各进程连接起来，把每个进程都当作一个远程提供者。跨进程的服务访问由一种保持接口不变的 RPC 机制来中介，使这种分布对消费者透明。一个需要留意之处是，跨进程调用会引入延迟，并且可能在调用中途失败，因此同步暴露它会阻塞调用方。因此，打算跨进程暴露的接口必须按异步契约来设计。

\subsection{6.3. 访问控制与沙箱隔离}\label{access-control-and-sandboxing}

对于由相互独立的组件组装而成的应用，保障其安全需要两种互补机制：(1) 约束组件可以访问哪些依赖，(2) 把不可信代码与宿主环境沙箱隔离。Cordis 通过依赖声明与拦截支持第一种；第二种则需要外部的沙箱。

基于能力的访问控制。依赖访问机制（第~5.1.4 节）已经构成了一种对经由代理中介的属性的访问控制：组件只能访问它已声明的依赖；未声明的访问会引发错误。这在结构上与基于能力的安全 {[}54--56{]} 相似，后者的权威来自对引用的持有，而非环境权威。inject 声明充当能力请求，而上下文代理充当能力中介。由于这些请求是静态声明的，组件所需的经由代理中介的完整能力集在其运行之前便已知晓，这让编排器可以在加载时审查并批准它们，而不是在访问发生时才发现。

这种中介通过拦截机制推广为细粒度策略。访问控制元数据可以由上下文携带，也可以由组件声明（定义~30），提供者在依赖被调用时查阅它，以决定某个请求是否被允许。例如，一个文件系统依赖可以携带元数据，声明组件可以读写哪些路径，提供者据此对每次调用进行检查。由于这种拦截存在于上下文之上，而非任何一方的代码之中，编排器可以调整它来约束任何组件对某个依赖的访问，而无需修改提供者，例如授予社区组件只读的数据库访问权限，而核心组件保留完全访问权限。此外，由于拦截只影响依赖被调用的方式，而不影响它是否被满足，因此可以在运行时安装、重新配置或移除它，而不会触发任何重载，也不会扰动依赖图。

69

沙箱隔离不可信组件。当组件代码不可信时，语言层面的访问控制是不够的，因为一个能访问宿主运行时的恶意组件可以直接触及底层对象，使这类检查形同虚设。沙箱隔离需要一个语言层面手段无法企及的执行边界，例如软件故障隔离 {[}57{]}、独立的语言运行时、沙箱化进程或虚拟化容器 {[}58{]}。无论采用何种机制，不可信组件都在它自己的沙箱化上下文中运行，并通过一座桥接访问宿主提供的依赖，这是对第~6.2 节跨进程调用的一般化：同样的透明性论证，使这种桥接访问与本地注入无从区分。在宿主一侧，这座桥接是一条普通纤维，其能力可以被上述访问控制所削减。

\subsection{6.4. 语言无关性与语言选择}\label{language-independence-and-selection}

尽管 Cordis 用 TypeScript 实现，但上下文范式与语言无关：时空可组合性只由它的两个可组合性维度来定义，因而可以在任何沿两个维度都满足一定要求的语言中实现。我们逐一分析这两个维度上的要求。

时间可组合性。在最基本的层面，时间可组合性要求闭包：一个可逆效应把操作与逆配对，而逆必须作为值被捕获，连同它所恢复的状态一起，以便在拆卸时重放。除此之外，组件的代码以及加载它所伴随的副作用，都必须能在运行时被引入和撤销。

语言如何满足这第二个要求，取决于它的执行模型。在托管运行时中，这表现为程序化的模块注册表：已加载的模块可以被从注册表中逐出，并在不再被引用后由垃圾回收回收；例如，Node.js 就暴露了这样一个注册表。\footnote{CommonJS 通过 require.cache 暴露模块缓存；ES 模块不提供公开的驱逐 API，不过模块仍可通过引擎内部接口来管理。}原生代码不暴露模块注册表，因此引入与撤销采取显式动态链接与解链的形式（例如 Unix 上的 dlopen/dlclose、Windows 上的 LoadLibrary/FreeLibrary）{[}59{]}，即把目标代码加载进运行中的进程，随后再把它卸载。WebAssembly 则视嵌入器而定选择两条路径之一：在托管嵌入器（例如 JavaScript 宿主）下，模块实例由宿主的回收器回收；在原生嵌入器（例如 Wasmtime）释放它时则被释放。在这些机制之上，可逆效应模型把加载当作上下文上的一个效应，其逆会撤销模块引入的符号、类型或处理器的注册。

空间可组合性。空间可组合性要求组件能够声明依赖、运行时能够提供并注入这些依赖的机制。这归结为一个依赖注入（DI）问题 {[}38{]}，它在两个随语言而异的层面上体现：依赖如何被类型化，以及依赖的访问如何被中介。

在类型层面，语言应当为开发者提供一种表达良类型化依赖访问的途径。消费者通过从上下文读取键来获得一个余效应，因此上下文类型（第~3.2.1 节）必须记录每个键的余效应。类型类（Haskell）{[}60{]} 与 trait（Rust）{[}61{]} 通过让提供者从它自己的模块出发、以 instance 或 impl {[}62{]} 扩展上下文类型来实现这一点。TypeScript 的模块扩充 {[}63{]} 同样让提供者模块把声明并入上下文类型。

在运行时层面，依赖访问必须被动态中介：键背后的余效应会随提供者的加载与卸载而改变，并且可能在不同的上下文间被不同地解析。因此，语言需要一种透明地介入访问、同时不改动消费者代码的途径，例如通过 JavaScript 的 Proxy 对象 {[}64{]}，或 Python 的描述符协议（\texttt{\_\_get\_\_}）{[}65{]}。若缺少这样的原语，运行时反射 {[}66, 67{]} 也能动态地中介访问，但代价是类型安全与开发者体验。

70

在这两个层面之上，元编程设施把类型化与中介一并提供。注解 {[}68{]} 与装饰器把元数据附着到声明上，处理器把它展开成中介访问的访问器；编译期元编程（例如 Rust 过程宏、Scala 宏 {[}69{]}、Zig comptime）为每个依赖发出一个带类型的声明以及这样一个访问器，从而免去通用的拦截原语。

\subsection{6.5. 相互依赖与组件粒度}\label{mutual-dependencies-and-component-granularity}

在响应式余效应模型中，依赖环只会让涉及的组件永久地保持不活跃：给定两个组件 \(A\) 和 \(B\)，若 \(A\) 需要一个由 \(B\) 提供的键，而 \(B\) 需要一个由 \(A\) 提供的键，那么两者的满足谓词都不可能成立。与并发系统中的死锁不同——死锁取决于调度，并且只能在发生时被检测到——这一条件仅凭依赖声明就可预测，因此运行时可以在组件加载时就报告它。

在实践中，大多数表面上的相互依赖都可以分解为更细粒度的组件，从而消除环。考虑两个组件：一个服务器（提供网络接口）和一个访问控制器（执行授权策略）。这两个组件双向交互：访问控制器中介到达服务器的请求，而服务器暴露一个用于修改访问控制策略的端点。单体设计会让两个组件互相依赖。然而，这两个交互方向是逻辑上相互独立的关切。把它们分解得到四个组件：server-core、accesscontrol-core、request-mediation（依赖两个核心，以便对到达的请求施加访问控制）与 policy-management（依赖两个核心，以便通过服务器暴露策略修改）。通过这种方式，环被消除了，因为两个核心互不依赖；只有集成组件同时依赖两者。

这种分解原则上总是可行，因为每个双向交互都可以分解为相互独立的单向绑定，但它会增加组件数量：在一般情况下，给定 \(n\) 个相互交互的组件，集成组件的数量可能随 \(n\) 二次增长，因为每对交互组件都可能需要为每个交互方向准备一个不同的组件。这不会影响正确性或运行时性能（组件是轻量的），而且更细的粒度可能是有益的：用户能够只加载他们需要的特定集成绑定，从而有效地提高系统的可组合性。不过，它可能影响开发者体验：更多的组件意味着更多的配置、更多的命名，以及在理解依赖图时更多的认知开销。

缓解这种粒度代价是一个工程关切，而非理论问题。实用的策略包括打包（即把相关细粒度组件归并为一个可安装单元）、基于约定的接线（即自动连接名字或类型匹配某个模式的组件），以及脚手架工具（即从声明式规范生成样板集成组件）。这些策略保留了无环模型的形式保证，同时把编写负担降低到更接近单体情形。

71

\subsection{6.6. 依赖类型化与版本管理}\label{dependency-typing-and-versioning}

在形式化模型中，依赖链接纯粹由键标识建立：提供键 \(k\) 的组件满足任何在其依赖集中声明 \(k\) 的组件。类型族 \(\nu _ { k }\) 保证单个编译单元内的类型层面一致，但当组件被独立开发和构建时——这是组件生态中的常见情形——这一保证会失效。这种失效导致两个截然不同的问题。

接口漂移。提供者可能在版本之间修改与 \(k\) 关联的接口（增加字段、改变方法签名、改变行为契约），而针对较早接口编译的消费者继续声明相同的键 \(k\)。依赖在余效应层面得到满足 \(( k \in \mathrm { d o m } ( \sigma ) )\)，但运行时的值不再符合消费者的预期，从而导致类型错误、方法未找到失败或静默的行为分歧 {[}70{]}。

键冲突。两个独立开发的提供者可能使用相同的键名 \(k\) 来表示完全无关的接口。由于单凭键标识就建立链接，期望某个提供者接口的消费者会不加任何兼容性检查地接受另一个提供者的值。与接口漂移不同——在那里提供者与消费者至少共享同一谱系——键冲突在预期类型与实际类型之间不存在任何关联，使得由此产生的失败难以预测、难以诊断。

两个问题都指向同一个缺口：余效应模型只提供名义链接（按键名），而不提供版本化或结构化链接（按接口兼容性）{[}71{]}。我们讨论填补这一缺口的三种方法，从与基础设施耦合最紧到最语言无关。

键命名空间化。把键空间从 \(K\) 扩展到 \(K \times P ,\)（其中 \(P\) 标识定义接口的包），从构造上消除了键冲突：本地名相同但独立开发的接口会占据不同的键。这是最直接的解决方案，但耦合也最紧：它把包命名空间嵌入形式化模型本身，使系统依赖外部的包注册表来提供键标识。

对等依赖。一种更松的耦合是通过宿主语言的包管理器来声明版本约束 {[}72{]}。这是 Cordis 目前采用的方法。组件依赖在语义上是对等依赖：组件并不在其内部捆绑依赖，而是期望运行时上下文提供它们。支持对等依赖的包管理器（例如，npm）可以强制版本兼容：如果提供某个键的包版本落在消费者声明的对等版本范围之外，不兼容性会在安装时被发现，而不是作为运行时失败浮现出来。不过，这种方法有两个局限：(1) 它依赖提供者忠实地遵循语义化版本，而这是一种无法强制执行的约定；(2) 包管理器通常把每个依赖解析为单一版本，这阻止了在一个应用内加载同一包多个版本的组件。

结构化兼容。一种完全语言无关的方法，会把成员判定 \(k \in \mathrm { d o m } ( \sigma )\) 替换为一个兼容性谓词，验证提供者的实际接口在结构上包含消费者的预期。这与结构化子类型 {[}73{]} 类似：当提供的接口是被要求的接口的子类型时，提供者满足消费者。挑战在于以语言无关的方式定义这个谓词：结构化兼容对记录类型（宽度子类型）是直接的，但对行为契约（例如前置/后置条件 {[}74{]}、效应规范 {[}22{]}）会变得复杂，而一旦参数多态引入有界量化 {[}75{]}，它便不可判定。

72

这三种方法处理问题的不同方面。设计一个统一依赖模型，把三种方法结合起来，同时保留余效应模型的动态组合保证，仍是一个开放问题。

\subsection{6.7. 与语言和操作系统的协同设计}\label{co-design-with-languages-and-operating-systems}

第~6.4 节指出了宿主语言为上下文范式所必须提供的最低限度。本节转向相反的问题：与范式协同设计的语言或操作系统，能在这一最低限度之上提供什么。

与语言的协同设计。围绕上下文范式设计的语言，可以在两个方面优于一个库：它赋予上下文的语义，以及它赋予效应与余效应的原语。

这样的语言可以在保持第~3.3 节上下文语义的同时，再次让上下文变为隐式的。命令式语言已经让每条语句针对一个隐式上下文运行，而那个单一上下文既不跟踪效应，也不解析余效应。上下文范式则区分多个上下文，其中一次操作要么修改它所运行的上下文，要么从中派生出另一个上下文（定义~27）。就地实现修改环境上下文，正如命令式语言所做的那样。派生实现则引入一个独立的上下文，语言必须为此提供一种构造。让上下文隐式化同时带来易用性与安全性的好处。(1) 在库实现中，每个涉及效应或余效应的函数都把上下文作为普通参数或接收者来取，如第~5.1 节所示。当语言隐式地提供上下文时，函数便不再需要取它。(2) 每个上下文都携带自己的生命周期状态与已提交视图（第~4.1 节）。库实现把上下文作为普通变量传递，因此组件可能通过闭包或全局变量，误触达另一个组件的上下文。它在那边安装的效应于是泄漏出自己的生命周期，而它在那边读取的余效应则逃逸出它的依赖规范。让上下文隐式化同时封堵了这两点。

这样的语言还可以让效应与余效应为其编译器所知。(1) 对效应而言，效应迭代器（定义~51）在每一步都分配一个闭包，以保存逆及它所恢复的状态。有了执行效应的语法，编译器就可以为整个迭代发出单个状态机，并把那些逆保存在它的帧中。(2) 对余效应而言，余效应规范可以被纳入类型系统，带来两个好处。其一，依赖环在编译期被报告，而不是留给运行时（第~6.5 节）。其二，依赖可以按其类型的结构而非仅按键标识来比较，正如行类型所做的那样 {[}28{]}，这是对第~6.6 节结构化兼容的类型层面支持。

与操作系统的协同设计。第~1.2.3 节观察到动态可组合性的一种粗粒度替代物，其中操作系统在进程粒度上提供时间可组合性，其上层的容器编排器在服务粒度上提供空间可组合性。与范式协同设计的操作系统将支持细粒度的组合，方法是把组件声明的余效应规范当作它所能触达的全部，并把自身资源作为余效应来提供。

这样的操作系统可以提供第~6.3 节推迟给语言之外机制的那个沙箱。它把组件限定在它声明的依赖之内，在组件加载时提供这些依赖，并使其内部不再有其他可触达之物，正如 WebAssembly 模块在实例化时从它的嵌入器接收导入那样 {[}76{]}。它还可以把第~3.2.3 节的余效应隔离与拦截作为自身的能力来提供，为每个组件以不同的方式绑定键，并中介它所提供的访问。

73

这样的操作系统还可以把自身资源作为余效应提供。位于边界之外的资源，在运行时按组件对每次获取记账的情况下是可逆的（第~6.1 节），而每个运行时都保有自己的记录。把资源作为余效应提供的操作系统只需保留这份记录一次，因为它是把资源分发出去的一方，可以把资源归属于提出请求的组件。内存与文件描述符是现成的候选，而为了恢复之目的跟踪它们，已在内核接口层面实现 {[}77, 78{]}。此外，操作系统可以让第~6.1 节只能扣留或补偿的部分操作变为可逆。以事务方式对持久化存储执行写入的系统可以回滚它 {[}79{]}，而构建在写时复制或不可变存储之上的系统通过移动一个指针即可到达较早的状态 {[}80, 81{]}。
